\begin{table}[htbp]
\centering
\caption{Cuadro 1: Programas y proyectos prioritarios de inversión}
\small Millones de pesos
\begin{tabular}{lccc}
\toprule
Concepto & PEF 2025 & PPEF 2026 & \% variación real \\
\midrule
1. Nuevos trenes & 77,655 & 104,576 & 34.7 \\
   Tren México – Querétaro & 31,062 & 10,492 & -66.2 \\
   Tren AIFA – Pachuca & 25,885 & 3,108 & -88.0 \\
   Tren Saltillo – Nuevo Laredo & 10,354 & 14,386 & 38.9 \\
   Tren Querétaro – Irapuato & 10,354 & 9,344 & -9.8 \\
   Material rodante & n.a. & 14,879 & n.a. \\
   Derecho de vía & n.a. & 21,432 & n.a. \\
   Tren México – SLP & n.a. & 9,076 & n.a. \\
   Tren Irapuato – Guadalajara & n.a. & 12,506 & n.a. \\
   Tren San Luis Potosí – Saltillo & n.a. & 9,353 & n.a. \\
2. Obras de Interconexión AIFA (Tren Lechería) & 277,334 & 600 & -99.8 \\
3. Tren México–Toluca & 10,354 & 7,408 & -28.5 \\
4. Trolebús Iztapaluca & 1,294 & 1,250 & -3.4 \\
5. Tren Maya (trenes, mantenimiento y tren Campeche) & 41,416 & 30,000 & -27.6 \\
6. Obras hidráulicas CONAGUA & 20,708 & 20,762 & 0.3 \\
7. Carreteras y Caminos & 10,354 & 27,720 & -62.6 \\
8. Defensa Nacional & 3,453 & 7,000 & -50.7 \\
9. Istmo de Tehuantepec & 25,885 & 25,000 & -3.4 \\
10. Marina & 2,283 & 2,000 & -12.4 \\
11. Inversión SSPC & 2,565 & 2,169 & -15.4 \\
12. PEMEX & 216,954 & 247,230 & 14.0 \\
13. CFE & 63,715 & 61,091 & -4.1 \\
Total & 476,638 & 536,806 & 11.2 \\
\midrule
\multicolumn{4}{l}{\footnotesize \textit{N.A. Indica que no se registró el proyecto para ese año.}}
\multicolumn{4}{l}{\footnotesize Fuente: Elaborado por el CIEP, con información de SHCP (2025a).}
\bottomrule
\end{tabular}
\end{table}
